%Some random things

\section{Sharper tails and restrictions on a possible limiting law - Random directions}
\label{sec:fraiture-distribution-addendum}

Write
\[
 h(m)=\frac{\sigma(m)}m,
 \qquad
 F_e(d)=\sum_{\substack{r\mid ed\\r\leq d}}r,
 \qquad
 g_e(m)=\sum_{\substack{r\mid m\\r\geq e}}\frac1r.
\]
Complementary divisors give the exact identity
\begin{equation}
 F_e(m/e)=m g_e(m)
 \qquad(e\mid m).
 \label{eq:addendum-reflection}
\end{equation}
Since $\Rcal=\N\setminus\{2,5\}$, retain the normalized empirical
probability measures from the preceding section and put
\[
 R(X)=\#(\Rcal\cap[1,X])=X+O(1),
 \qquad
 \mu_X=\frac1{R(X)}
       \sum_{\substack{N\leq X\\N\in\Rcal}}
          \delta_{f(N)/N},
 \qquad
 G_E=\bigcup_{1\leq e\leq E}F_e(\N).
\]
Replacing $R(X)^{-1}$ by $X^{-1}$ in any bounded empirical count
changes it by $O(X^{-1})$.
As in the preceding sections, $\log_j$ denotes the $j$-fold iterated
natural logarithm.  The existing large-cofactor moment estimate gives,
for every fixed integer $k\geq2$,
\begin{equation}
 \sum_{m\leq Z}\ \sum_{\substack{e\mid m\\e>E}}g_e(m)^k
 \leq C_kZE^{1-k}
 \qquad(E,Z\geq1).
 \label{eq:addendum-large-cofactor-moment}
\end{equation}
In particular, by Markov's inequality,
\begin{equation}
 \#\left\{(m,e):
  \begin{array}{c}
   e\mid m,\ e>E,\ F_e(m/e)\leq X,\\
   m\leq tF_e(m/e)
  \end{array}\right\}
 \leq C_kt^{k+1}E^{1-k}X.
 \label{eq:addendum-large-cofactor-witnesses}
\end{equation}
The constants are allowed to depend on $k$, and all uses below respect
the order in which the parameters are fixed.

% Quantitative addendum for the original divisor-prefix minimum.
% This fragment uses the notation and theorem environments of the main paper.
% Its only external cross-references are lem:sigma-rate, lem:moment, and
% thm:fraiture-representability.

\subsection{Two preliminary rough-kernel estimates}

For a real number $y\geq2$, write
\[
 Q_y=\prod_{p\leq y}p,
 \qquad
 \Delta_y=\prod_{p\leq y}\left(1-\frac1p\right),
 \qquad
 M_y=\operatorname{lcm}(1,\ldots,\lfloor y\rfloor).
\]

\begin{lemma}[Fixed-modulus divisor-sum normality]
\label{lem:qadd:fixed-modulus-normality}
For every fixed positive integer $V$,
\[
 \#\{m\leq Y:V\nmid\sigma(m)\}=o_V(Y)
 \qquad(Y\to\infty).
\]
\end{lemma}

\begin{proof}
It suffices to prove the assertion when $V=q$ is a prime power.
For every prime $r\equiv-1\pmod q$, the condition $r\Vert m$
implies
\[
 q\mid r+1=\sigma(r)\mid\sigma(m).
\]
Therefore, if $q\nmid\sigma(m)$, none of these primes divides
$m$ exactly once. For any finite set $\mathcal P$ of such primes,
the Chinese remainder theorem gives
\[
 \limsup_{Y\to\infty}\frac1Y
       \#\{m\leq Y:q\nmid\sigma(m)\}
 \leq\prod_{r\in\mathcal P}
            \left(1-\frac1r+\frac1{r^2}\right).
\]
The sum of the reciprocals of the primes in the reduced residue
class $-1\pmod q$ diverges. Letting $\mathcal P$ increase through
this class makes the displayed product tend to zero. Finally,
factor a general fixed $V$ into its finitely many prime-power
divisors and apply the union bound.
\end{proof}

\begin{lemma}[Two saturated-kernel tails]
\label{lem:qadd:kernel-tails}
For $y,H\geq2$, put
\[
 S(y,H)=\sum_{\substack{L\mid M_y\\L>H}}
                  \frac{\tau(L)}L.
\]
If $E\to\infty$ and
\[
 P=\left\lfloor
 E^{\,1+(2+(\log_3 E)^{-1/2})\log_2E/\log_3E}
 \right\rfloor,
\]
then
\begin{equation}
\label{eq:qadd:fixed-kernel-tail}
 \log S(E,P/E)
 \leq-\frac{\log_2E}{\sqrt{\log_3E}}
      +O\left(\frac{\log_2E}{\log_3E}\right).
\end{equation}
If instead $J\to\infty$,
\[
 W=\sqrt{\frac{J\log J}{2}},
 \qquad F=\lfloor e^W\rfloor,
\]
then, for every fixed $C>0$,
\begin{equation}
\label{eq:qadd:moving-kernel-tail}
 S\left(F,\frac{e^J}{CW}\right)
 \leq\exp\{-W+o(W)\}.
\end{equation}
\end{lemma}

\begin{proof}
For $0<\alpha\leq1/4$, Rankin's inequality and the Euler product
for the divisor function give
\begin{align}
 S(y,H)
 &\leq H^{-\alpha}
       \prod_{p\leq y}\sum_{a\geq0}
                \frac{a+1}{p^{a(1-\alpha)}}
  =H^{-\alpha}
       \prod_{p\leq y}(1-p^{-1+\alpha})^{-2}.
 \label{eq:qadd:rankin-product}
\end{align}
Mertens' estimate for $\sum_{p\leq y}1/p$, Chebyshev's bound,
and partial summation show that, when
$z=\alpha\log y\to\infty$,
\begin{align}
 \sum_{p\leq y}p^{-1+\alpha}
 &=\log_2y+
   O\left(\int_0^z\frac{e^u-1}{u}\,du\right)
  =\log_2y+O\left(\frac{e^z}{z}\right).
 \label{eq:qadd:sharp-prime-sum}
\end{align}
The higher terms in the logarithm of the Euler product
in~\eqref{eq:qadd:rankin-product} are $O(1)$ uniformly for
$\alpha\leq1/4$.

For~\eqref{eq:qadd:fixed-kernel-tail}, take
$\alpha=\log_3E/\log E$. Then
\[
 \alpha\log(P/E)
 =\left(2+\frac1{\sqrt{\log_3E}}+o(1)\right)\log_2E,
\]
whereas~\eqref{eq:qadd:sharp-prime-sum} gives
\[
 2\sum_{p\leq E}p^{-1+\alpha}
 =2\log_2E+
   O\left(\frac{\log_2E}{\log_3E}\right).
\]
Substitution in~\eqref{eq:qadd:rankin-product} proves the claim.

For~\eqref{eq:qadd:moving-kernel-tail}, take $\alpha=W/J$.
Since
\[
 \alpha\log F=\frac12\log J+o(1),
 \qquad
 \alpha\log\left(\frac{e^J}{CW}\right)=W+o(W),
\]
equation~\eqref{eq:qadd:sharp-prime-sum} gives
\[
 \sum_{p\leq F}p^{-1+\alpha}
 \ll\log W+\frac{\sqrt J}{\log J}=o(W).
\]
The second assertion follows once more
from~\eqref{eq:qadd:rankin-product}.
\end{proof}

\subsection{An almost-logarithmic fixed-level upper tail}

\begin{theorem}[Squarefree fixed levels]
\label{thm:qadd:almost-log-tail}
As $T\to\infty$,
\begin{equation}
\label{eq:qadd:almost-log-tail}
 \liminf_{X\to\infty}\frac1X
 \#\left\{N\leq X:
       N\in\mathcal R,\ N\text{ squarefree},\ f(N)>TN
       \right\}
 \geq\left(\frac{e^{-\gamma}}2-o(1)\right)
       \frac{\log_3T}{(\log T)(\log_2T)}.
\end{equation}
In particular, the same lower bound holds without the
squarefreeness restriction.
\end{theorem}

\begin{proof}
Fix an integer $j\geq3$ and put
\begin{align*}
 b_j&=\frac{j+1}{j-1},\\
 E&=\left\lceil T^{b_j}
                (\log T)^{3/(j-1)}\right\rceil,\\
 P&=\left\lfloor
 E^{\,1+(2+(\log_3E)^{-1/2})\log_2E/\log_3E}
 \right\rfloor.
\end{align*}
All parameters in this proof are fixed while $X\to\infty$.
For $2\leq e\leq E$ and $L\mid M_E$, define
\[
 C_e(L)=\sum_{\substack{r\mid L\\r<e}}\frac Lr,
 \qquad
 k_e(L)=\frac{C_e(L)}{\operatorname{sm}_P(C_e(L))},
\]
where $\operatorname{sm}_P(v)$ is the largest divisor of $v$
all of whose prime factors are at most $P$. Set
\[
 K=\operatorname{lcm}_{\substack{2\leq e\leq E\\L\mid M_E}}
                     k_e(L),
 \qquad V=KQ_P.
\]
Lemma~\ref{lem:qadd:fixed-modulus-normality} shows that the
sources $m\leq EX$ with $V\nmid\sigma(m)$ number $o_{E,P}(X)$.
Since each such source has at most $E$ possible cofactors not
exceeding $E$, their entire bounded-cofactor image also has
cardinality $o_{E,P}(X)$.

Let $N\leq X$ be $P$-rough and suppose that
$N=F_e(m/e)$ with $e\leq E$ and $V\mid\sigma(m)$.
Here $m\leq eN\leq EX$. The case $e=1$ is impossible, since
it would give $N=\sigma(m)$ although $Q_P\mid\sigma(m)$.
For $e\geq2$, put
\[
 L=\operatorname{lcm}\{r:r\mid m,\ r<e\},
 \qquad C=C_e(L),
 \qquad q=\sigma(m)-N.
\]
Reflection of the divisor list gives the exact identity
\begin{equation}
\label{eq:qadd:complementary-kernel}
 q=\frac mL C.
\end{equation}
Since $Q_P\mid\sigma(m)$ and $(N,Q_P)=1$, the integer $q$ is
$P$-rough. Consequently both $u=m/L$ and $C$ are $P$-rough.
For every prime $p\mid e$, one has $p\leq E<P$ and therefore
$v_p(L)=v_p(m)\geq v_p(e)$, because $m=Lu$ and $p\nmid u$.
Thus $e\mid L$. Moreover $L\mid M_E$ and
\[
 C\geq L\geq e\geq2,
 \qquad C\leq EL.
\]
Thus $C>P$, $L>P/E$, and $k_e(L)=C$. By the definition of
$K$, we have $C\mid K\mid\sigma(m)$, while
\eqref{eq:qadd:complementary-kernel} gives $C\mid q$.
Therefore $C\mid N$.

For fixed $e,L$, the $P$-rough multiples of $C$ have natural
density $\Delta_P/C$. Since $C\geq L$ and $e\mid L$, summing
over the finitely many admissible pairs gives total upper density
at most
\[
 \Delta_P\sum_{\substack{L\mid M_E\\L>P/E}}
                         \frac{\tau(L)}L
 =\Delta_P S(E,P/E)
 =o(\Delta_P)
 \qquad(T\to\infty).
\]
By Lemma~\ref{lem:moment}, the targets admitting a
representation with $e>E$ and $m\leq TN$ have upper density
at most
\[
 C_jT^{j+1}E^{1-j}
 \leq\frac{C_j}{(\log T)^3}
 =o(\Delta_P).
\]
Finally, a $P$-rough target which is not squarefree is divisible
by $p^2$ for some $p>P$, so these targets have upper density
at most
\[
 \sum_{p>P}\frac1{p^2}\ll\frac1P=o(\Delta_P).
\]
Theorem~\ref{thm:fraiture-representability} discards at most
two additional targets. Subtracting all three exceptional sets
from the $P$-rough cylinder yields
\[
 \liminf_{X\to\infty}\frac1X
 \#\{N\leq X:N\in\mathcal R,\ N\text{ squarefree},\ f(N)>TN\}
 \geq(1-o(1))\Delta_P.
\]
Mertens' theorem and the definitions of $P,E$ give
\[
 \Delta_P
 =\left(\frac{e^{-\gamma}}{2b_j}+o(1)\right)
       \frac{\log_3T}{(\log T)(\log_2T)}.
\]
First let $T\to\infty$ for each fixed $j$, and then let
$j\to\infty$. Since $b_j\to1$, this proves
\eqref{eq:qadd:almost-log-tail}.
\end{proof}

\subsection{Subexponential growing levels on squarefree targets}

\begin{theorem}[Subexponential squarefree growing levels]
\label{thm:qadd:growing-level}
For every fixed $\eta>0$ there exists $c_\eta>0$ such that,
for all sufficiently large $X$,
\begin{equation}
\label{eq:qadd:growing-level}
 \#\left\{\frac X2<N\leq X:
       N\in\mathcal R,\ N\text{ squarefree},\quad
       \frac{f(N)}N>
       \exp\left[
         \left(\frac1{\sqrt2}-\eta\right)
         \sqrt{\log_3N\,\log_4N}
       \right]\right\}
 \geq c_\eta\frac X{\log_3X}.
\end{equation}
Consequently, for every fixed $s>0$ and $\eta>0$,
\begin{equation}
\label{eq:qadd:positive-moment-growth}
 \frac1X\sum_{\substack{N\leq X\\N\in\mathcal R}}
       \left(\frac{f(N)}N\right)^s
 \geq
 \exp\left[
       \left(\frac{s}{\sqrt2}-\eta\right)
       \sqrt{\log_3X\,\log_4X}
      \right]
\end{equation}
for all sufficiently large $X$.
\end{theorem}

\begin{proof}
Fix $0<\varepsilon<1$ sufficiently small in terms of $\eta$.
Choose a fixed $\rho>0$ sufficiently small, depending only on
the absolute constants in Lemma~\ref{lem:sigma-rate}, and put
\[
 P=\left\lfloor\rho\frac{\log_2X}{\log_3X}\right\rfloor,
 \qquad J=\log P,
 \qquad W=\sqrt{\frac{J\log J}{2}},
 \qquad F=\lfloor e^W\rfloor,
 \qquad t=e^{(1-\varepsilon)W}.
\]
In particular $F<P$ for all sufficiently large $X$.
Complete periods modulo $Q_P$ and Mertens' theorem give
\begin{equation}
\label{eq:qadd:rough-target-count}
 \#\{X/2<N\leq X:(N,Q_P)=1\}
 =\frac X2\Delta_P+O(Q_P)
 \asymp\frac X{\log_3X},
\end{equation}
because $Q_P=\exp(O(P))=X^{o(1)}$ and
$\Delta_P\asymp1/J$. The nonsquarefree members of this interval
number at most
\[
 X\sum_{p>P}\frac1{p^2}\ll\frac XP
 =o(X\Delta_P).
\]

Write $Y=tX$. Lemma~\ref{lem:sigma-rate} and the union bound
show that the number of sources $m\leq Y$ for which some prime
$p\leq P$ fails to divide $\sigma(m)$ is at most
\begin{align}
 O(\sqrt Y)
 +O\left(
      Y\pi(P)
       \exp\left[-c_1\frac{\log_2Y}{P}\right]
    \right).
 \label{eq:qadd:bad-source-count}
\end{align}
The first term treats $p=2$. The chosen range of $P$ is
permitted by Lemma~\ref{lem:sigma-rate} after decreasing
$\rho$ if necessary. Since
\[
 Ft=e^{(2-\varepsilon)W}
      =(\log_2X)^{o(1)},
 \qquad
 \exp\left[-c_1\frac{\log_2Y}{P}\right]
      =(\log_2X)^{-c_1/\rho+o(1)},
\]
we may also ensure $c_1/\rho>3$. Multiplying
\eqref{eq:qadd:bad-source-count} by the at most $F$ possible
cofactors per source then gives $o(X\Delta_P)$ target--cofactor
pairs.

Consider a remaining $P$-rough target $N$ admitting a witness
$N=F_e(m/e)$ with $m\leq tN$ and $e\leq F$.
Again $e=1$ is impossible, since it would give
$N=\sigma(m)$ although $2\mid\sigma(m)$.
For $e\geq2$, define $L,C,q$ exactly as
in~\eqref{eq:qadd:complementary-kernel}. Since every prime
$p\leq P$ divides $\sigma(m)$ but not $N$, the integer
$q=(m/L)C$ is $P$-rough. Writing $u=m/L$, both $u$ and $C$
are therefore $P$-rough. Since $e\leq F<P$ and $e\mid m=Lu$,
every prime-power divisor of $e$ already divides $L$.
Consequently
\[
 m=Lu,
 \qquad (u,Q_P)=1,
 \qquad (C,Q_P)=1,
 \qquad e\mid L\mid M_F.
\]
As before $C\geq L\geq2$, so $C>P$. On the other hand,
\[
 C\leq\sigma(L)
 \leq L\prod_{p\leq F}\left(1-\frac1p\right)^{-1}
 \ll L\log F.
\]
There is therefore an absolute constant $C_0>0$ such that
\begin{equation}
\label{eq:qadd:moving-kernel-lower-bound}
 L>\frac P{C_0W}.
\end{equation}

Uniformly for all $L\mid M_F$,
\[
 \log L\leq\log M_F\ll F=o(\log X),
 \qquad
 \log Q_P\ll P=o(\log X).
\]
It follows that $tX/L\geq Q_P$ for all sufficiently large $X$.
Complete periods modulo $Q_P$ consequently give
\[
 \#\{u\leq tX/L:(u,Q_P)=1\}
 \leq\frac{2tX\Delta_P}{L}.
\]
For each fixed $L$, the admissible cofactors satisfy $e\mid L$,
so there are at most $\tau(L)$ possibilities. Hence
\eqref{eq:qadd:moving-kernel-lower-bound} and
Lemma~\ref{lem:qadd:kernel-tails} bound the entire remaining
low-cofactor witness count by
\begin{align*}
 2tX\Delta_P
     \sum_{\substack{L\mid M_F\\L>P/(C_0W)}}
                \frac{\tau(L)}L
 &\leq X\Delta_P
      \exp\{-\varepsilon W+o(W)\}\\
 &=o(X\Delta_P).
\end{align*}

Choose a fixed integer $k\geq2$ so large that
$(k+1)\varepsilon>2$. Lemma~\ref{lem:moment} bounds the
targets with a witness $m\leq tN$ and cofactor $e>F$ by
\[
 C_kt^{k+1}F^{1-k}X.
\]
After division by $X\Delta_P$, this is at most
\[
 O_k\left(
 J\exp\{[2-(k+1)\varepsilon+o(1)]W\}
 \right)=o(1).
\]
Thus a positive proportion of the squarefree targets
in~\eqref{eq:qadd:rough-target-count} has no representation
with $m\leq tN$. By
Theorem~\ref{thm:fraiture-representability}, all but at most
two of these targets are represented. Finally,
\[
 \log t
 =\left(\frac{1-\varepsilon}{\sqrt2}+o(1)\right)
          \sqrt{\log_3X\,\log_4X}.
\]
Choose $\varepsilon$ sufficiently small in terms of $\eta$.
For $X/2<N\leq X$, replacing $X$ by $N$ in these iterated
logarithms changes the right-hand side by $o(1)$ relative to
its main term. This proves~\eqref{eq:qadd:growing-level}.

For~\eqref{eq:qadd:positive-moment-growth}, apply the first
assertion with a sufficiently small auxiliary value of $\eta$,
raise the resulting ratios to the power $s$, and note that
\[
 \log\log_3X=\log_4X
 =o\left(\sqrt{\log_3X\,\log_4X}\right).
\]
The factor $c_\eta/\log_3X$ in the proportion of surviving
targets can therefore be absorbed into the prescribed exponent
loss $\eta$.
\end{proof}

% Distributional addendum for Erdos problem 1054.
% Input this fragment after the quantitative fixed-level theorem.
% Assumes the collaborative manuscript's definitions of F_e, g_e, h,
% R(X), mu_X, G_E, and its large-cofactor moment lemma.
% All new labels have the dadd namespace; external references are
% supplied in footnotes so that no bibliography changes are needed.

\subsection{The exact meaning of the bounded-cofactor conjecture}

For a set of positive integers $A$, write
\[
 \overline d(A)=\limsup_{X\to\infty}
       \frac{\#(A\cap[1,X])}{X},
 \qquad
 \underline d(A)=\liminf_{X\to\infty}
       \frac{\#(A\cap[1,X])}{X}.
\]

\begin{proposition}[Bounded cofactors are equivalent to tightness]
\label{prop:dadd:tightness}
The empirical probability measures $\mu_X$ are tight on
$[0,\infty)$ if and only if
\begin{equation}
\label{eq:dadd:tightness-criterion}
 \lim_{E\to\infty}
       \overline d(\mathbb N\setminus G_E)=0.
\end{equation}
Equivalently, the lower asymptotic densities of the finite
bounded-cofactor unions tend to one:
\[
 \lim_{E\to\infty}\underline d(G_E)=1.
\]
No asymptotic density is asserted for any individual $G_E$.
\end{proposition}

\begin{proof}
The summand $d$ occurs in $F_e(d)$, so $F_e(d)\geq d$. Hence
\[
 N\in G_E
 \quad\Longrightarrow\quad
 f(N)\leq e d\leq EN
 \quad\text{for some }e\leq E.
\]
Consequently
\begin{equation}
\label{eq:dadd:tightness-easy}
 \limsup_{X\to\infty}\mu_X((E,\infty])
 \leq\overline d(\mathbb N\setminus G_E).
\end{equation}
Here and below $R(X)=X+O(1)$ permits either normalization.
Thus~\eqref{eq:dadd:tightness-criterion} implies tightness.

Conversely, fix $A>0$ and an integer $k\geq2$. Every represented
$N\notin G_E$ with $f(N)\leq AN$ has a witnessing cofactor
$e>E$. The large-cofactor estimate therefore gives
\begin{equation}
\label{eq:dadd:tightness-reverse}
 \overline d(\mathbb N\setminus G_E)
 \leq\limsup_{X\to\infty}\mu_X((A,\infty])
      +C_kA^{k+1}E^{1-k}.
\end{equation}
First let $E\to\infty$, then let $A\to\infty$. Tightness makes
the remaining upper limit vanish. Finally,
$\overline d(\mathbb N\setminus G_E)=1-\underline d(G_E)$
identically, without assuming that $G_E$ has a density.
\end{proof}

\begin{corollary}[Every fixed cofactor union has a genuine defect]
\label{cor:dadd:cofactor-defect}
For every fixed $E\geq1$,
$\underline d(\mathbb N\setminus G_E)>0$. More precisely,
as $E\to\infty$,
\begin{equation}
\label{eq:dadd:cofactor-defect}
 \underline d\{N:N\text{ squarefree},\ N\notin G_E\}
 \geq
 \left(\frac{e^{-\gamma}}2-o(1)\right)
 \frac{\log_3E}{(\log E)(\log_2E)}.
\end{equation}
\end{corollary}

\begin{proof}
The argument preceding~\eqref{eq:dadd:tightness-easy} gives
\[
 \{N\in\mathcal R:N\text{ squarefree},\ f(N)>EN\}
 \subseteq\mathbb N\setminus G_E.
\]
Apply Theorem~\ref{thm:qadd:almost-log-tail} with $T=E$.
For a fixed smaller $E$, use any sufficiently large fixed $T\geq E$
instead; its fixed-level lower bound is strictly positive.
\end{proof}

\subsection{The limiting witness measure}

For $t>0$ and $N\geq1$, define the complete and proper witness
multiplicities by
\begin{align*}
 r_t(N)
 &=\#\{(e,d):F_e(d)=N,\ ed\leq tN\},\\
 r_t^*(N)
 &=\#\{(e,d):e\geq2,\ F_e(d)=N,\ ed\leq tN\}.
\end{align*}
Put $L_e=\operatorname{lcm}(1,\ldots,e)$, and, for a residue class
$a\pmod{L_e}$, let
\[
 c_e(a)=\sum_{\substack{1\leq j<e\\j\mid a}}\frac1j.
\]
Write $\mathcal D$ for the Davenport probability law of
$h(m)=\sigma(m)/m$. The Davenport theorem on arithmetic
progressions\footnote{P.~Pollack, \emph{Palindromic sums of proper
divisors}, Integers \textbf{15A} (2015), article A13, Lemma~2.}
gives an atomless measure $\nu_{a,e}$ characterized by
\[
 \nu_{a,e}(B)
 =\lim_{Y\to\infty}\frac1Y
    \#\{m\leq Y:m\equiv a\pmod{L_e},\ h(m)\in B\}
\]
for every continuity set $B$. In particular,
\begin{equation}
\label{eq:dadd:progression-domination}
 \sum_{a\bmod L_e}\nu_{a,e}=\mathcal D,
 \qquad 0\leq\nu_{a,e}\leq\mathcal D.
\end{equation}

\begin{proposition}[Explicit convergent witness intensities]
\label{prop:dadd:witness-means}
For each $t>0$ the limits
\[
 W(t)=\lim_{X\to\infty}\frac1X\sum_{N\leq X}r_t(N),
 \qquad
 W^*(t)=\lim_{X\to\infty}\frac1X\sum_{N\leq X}r_t^*(N)
\]
exist, are finite, and are continuous. More explicitly,
\begin{align}
 W(t)&=\sum_{e\geq1}w_e(t),
 &W^*(t)&=\sum_{e\geq2}w_e(t),
 \label{eq:dadd:witness-sums}\\
 w_e(t)
 &=\sum_{\substack{a\bmod L_e\\e\mid a}}
      \int_{h-c_e(a)\geq1/t}
       \frac{d\nu_{a,e}(h)}{h-c_e(a)}.
 \label{eq:dadd:witness-integral}
\end{align}
Both series converge uniformly for $0<t\leq T$, for every fixed
$T>0$, and $W(t)\to0$ as $t\downarrow0$.
\end{proposition}

\begin{proof}
For $m\equiv a\pmod{L_e}$ with $e\mid a$, complementary divisors
give the exact identities
\[
 F_e(m/e)=m g_e(m),
 \qquad
 g_e(m)=h(m)-c_e(a).
\]
In particular $g_e(m)\geq1/e$. A witness with target at most $X$
and ratio at most $t$ is characterized by
\[
 g_e(m)\geq1/t,
 \qquad m g_e(m)\leq X.
\]
Partition the variable $u=m/X\in(0,t]$ into finitely many
intervals, apply the progression-restricted Davenport theorem on
each interval, and then let the mesh tend to zero. Atomlessness
removes every boundary contribution. Tonelli's theorem gives
\begin{align*}
 &\lim_{X\to\infty}\frac1X
   \#\{m\equiv a\pmod{L_e}:
       g_e(m)\geq1/t,\ m g_e(m)\leq X\}\\
 &\quad=
 \int_0^t\nu_{a,e}
       \{h:h-c_e(a)\geq1/t,\ h-c_e(a)\leq1/u\}\,du\\
 &\quad=
 \int_{h-c_e(a)\geq1/t}
       \frac{d\nu_{a,e}(h)}{h-c_e(a)}.
\end{align*}
For the cofactor tail, the existing large-cofactor moment lemma
counts witnessing pairs themselves:
\begin{align}
 \frac1X\sum_{N\leq X}
  \#\{(e,d):e>E,\ F_e(d)=N,\ ed\leq tN\}
 &\leq\frac{t^k}{X}
   \sum_{m\leq tX}\sum_{\substack{e\mid m\\e>E}}g_e(m)^k
 \nonumber\\
 &\leq C_kt^{k+1}E^{1-k}.
 \label{eq:dadd:witness-tail}
\end{align}
This proves both limits and uniform convergence on $0<t\leq T$.
Each fixed $w_e(t)$ is continuous by atomlessness, and tends to
zero with $t$ since its integrand is bounded by $t$ on its domain.
The uniform tail proves all remaining assertions.
\end{proof}

\subsection{Every subsequential finite law is singular continuous}

\begin{theorem}[One universal singular carrier]
\label{thm:dadd:universal-singularity}
There is a fixed $\sigma$-compact, Lebesgue-null set
$\mathcal Z\subset(0,\infty)$ such that every weak subsequential
limit $\mu$ of $\mu_X$ on $[0,\infty]$ satisfies
\[
 \mu(\{0\})=0,
 \qquad
 \mu\bigl((0,\infty)\setminus\mathcal Z\bigr)=0,
 \qquad
 \mu(\{x\})=0\quad(0<x<\infty).
\]
Hence its restriction to $[0,\infty)$ is purely singular
continuous. If~\eqref{eq:dadd:tightness-criterion} holds, every
subsequential limit is a singular continuous probability law on
$[0,\infty)$; in particular, none admits a Lebesgue density.
Neither uniqueness of the limit nor tightness is asserted
unconditionally.
\end{theorem}

\begin{proof}
Erd\H{o}s proved that the continuous Davenport law $\mathcal D$ is
purely Lebesgue singular.\footnote{P.~Erd\H{o}s, \emph{On the
distribution of numbers of the form $\sigma(n)/n$ and on some
related questions}, Pacific J. Math. \textbf{52} (1974), 59--65.}
By regularity there is a $\sigma$-compact Lebesgue-null set
$K\subset[1,\infty)$ with $\mathcal D(K)=1$.
For $c\in\mathbb R$, write
\[
 \psi_c(h)=\frac1{h-c}\qquad(h>c),
\]
and define the positive witness-birth measure
\begin{equation}
\label{eq:dadd:singular-envelope}
 \zeta
 =\sum_{e\geq1}
   \sum_{\substack{a\bmod L_e\\e\mid a}}
    (\psi_{c_e(a)})_*
       \bigl(\boldsymbol1_{\{h>c_e(a)\}}
               \psi_{c_e(a)}(h)\,d\nu_{a,e}(h)\bigr).
\end{equation}
Proposition~\ref{prop:dadd:witness-means} gives
\[
 \zeta((0,t])=W(t)<\infty\qquad(t>0),
\]
so $\zeta$ is locally finite and atomless. By
\eqref{eq:dadd:progression-domination}, every $\nu_{a,e}$ is
carried by $K$. Therefore $\zeta$ is carried by
\begin{equation}
\label{eq:dadd:universal-null-carrier}
 \mathcal Z
 =\bigcup_{e\geq1}
    \bigcup_{\substack{a\bmod L_e\\e\mid a}}
    \bigcup_{j\geq1}
      \psi_{c_e(a)}
       \bigl(K\cap[c_e(a)+1/j,j]\bigr).
\end{equation}
Each map in the last display is Lipschitz on its indicated set.
Thus $\mathcal Z$ is $\sigma$-compact and Lebesgue null.

At finite scale, define the empirical witness-birth measure
\[
 \zeta_X
 =\frac1X\sum_{N\leq X}
    \sum_{\substack{e,d\geq1\\F_e(d)=N}}
       \delta_{ed/N}.
\]
On every bounded ratio interval it is finite. Proposition
\ref{prop:dadd:witness-means} and continuity of $W$ show that
$\zeta_X\to\zeta$ vaguely on $(0,\infty)$. For every
nonnegative $\varphi\in C_c((0,\infty))$, choosing a minimizing
witness separately for each represented target gives
\[
 \int\varphi\,d\mu_X
 \leq\frac X{R(X)}\int\varphi\,d\zeta_X.
\]
Along any weakly convergent subsequence $\mu_{X_j}\Rightarrow\mu$,
pass to the limit to obtain
\[
 \int\varphi\,d\mu\leq\int\varphi\,d\zeta.
\]
Hence $\mu|_{(0,\infty)}\leq\zeta$. This proves both singularity
on the common carrier and the absence of positive finite atoms.

Finally, the same witness domination and portmanteau give, for
every $t>0$,
\[
 \mu(\{0\})
 \leq\mu([0,t))
 \leq\liminf_{j\to\infty}\mu_{X_j}([0,t))
 \leq W(t).
\]
Let $t\downarrow0$ and use
Proposition~\ref{prop:dadd:witness-means}. Tightness is exactly
the exclusion of an atom at $\infty$, so the final assertion
follows from Proposition~\ref{prop:dadd:tightness}.
\end{proof}

\begin{corollary}[Heavy tails and divergent arithmetic and geometric means]
\label{cor:dadd:subsequential-tails}
Every compactified subsequential law $\mu$ satisfies, as
$T\to\infty$,
\begin{equation}
\label{eq:dadd:subsequential-tail}
 \mu((T,\infty])
 \geq\left(\frac{e^{-\gamma}}2-o(1)\right)
       \frac{\log_3T}{(\log T)(\log_2T)}.
\end{equation}
Moreover, for every $s>0$,
\[
 \lim_{X\to\infty}\int x^s\,d\mu_X(x)=\infty,
 \qquad
 \int_{[0,\infty]}x^s\,d\mu(x)=\infty.
\]
The logarithmic means also satisfy
\[
 \lim_{X\to\infty}\int\log x\,d\mu_X(x)=\infty,
 \qquad
 \int_{[0,\infty]}\log x\,d\mu(x)=\infty.
\]
In particular, the geometric means of $f(N)/N$ tend to infinity;
even assuming tightness, every subsequential probability law has
infinite logarithmic mean and infinite moments of all positive orders.
\end{corollary}

\begin{proof}
The closed-set direction of portmanteau and
Theorem~\ref{thm:qadd:almost-log-tail} give
\[
 \mu([T,\infty])
 \geq\limsup_{j\to\infty}\mu_{X_j}([T,\infty])
 \geq\liminf_{X\to\infty}\mu_X((T,\infty]).
\]
Theorem~\ref{thm:dadd:universal-singularity} excludes an atom at
$T$, proving~\eqref{eq:dadd:subsequential-tail}. For either $\mu_X$
or $\mu$, the $s$th moment is at least $T^s$ times its mass on
$(T,\infty]$. First take the lower limit in $X$ when appropriate,
and then let $T\to\infty$ in the preceding lower bound.

For the logarithmic means, layer cake and Fatou give
\begin{align*}
 \liminf_{X\to\infty}\int\log^+x\,d\mu_X(x)
 &\geq\int_{T_0}^B
       \liminf_{X\to\infty}\mu_X((T,\infty])\,\frac{dT}{T}\\
 &\geq\left(\frac{e^{-\gamma}}4-o(1)\right)
           (\log_3B)^2
 \qquad(B\to\infty).
\end{align*}
The same calculation using~\eqref{eq:dadd:subsequential-tail}
shows $\int\log^+x\,d\mu(x)=\infty$. On the other hand, the
existing small-value estimate, Theorem~\ref{thm:density}, gives
for some $c>0$ and all sufficiently large $u$
\[
 \sup_X\mu_X([0,e^{-u}])
 \ll\exp\bigl(-\exp(e^{cu})\bigr).
\]
Layer cake therefore bounds $\int\log^-x\,d\mu_X(x)$ uniformly
in $X$; portmanteau gives the same bound for $\mu$.
Subtracting the bounded negative parts proves both logarithmic
assertions.
\end{proof}

\subsection{The exact collision obstruction}

\begin{lemma}[The range of the full divisor sum has density zero]
\label{lem:dadd:sigma-range}
The set $\{\sigma(m):m\geq1\}$ has asymptotic density zero.
\end{lemma}

\begin{proof}
If $N=\sigma(m)\leq X$, then $m\leq X$. For every fixed $q$,
the values divisible by $q$ number at most $X/q$, while all
remaining values arise from $o_q(X)$ exceptional sources by
Lemma~\ref{lem:qadd:fixed-modulus-normality}. Thus
\[
 \overline d(\operatorname{range}(\sigma))\leq\frac1q.
\]
Letting $q\to\infty$ proves the claim.
\end{proof}

\begin{proposition}[Convergence is exactly a collision question]
\label{prop:dadd:collision-criterion}
Put
\[
 K_X(t)=\frac1X\sum_{N\leq X}(r_t^*(N)-1)_+.
\]
For every fixed $t>0$,
\begin{equation}
\label{eq:dadd:collision-criterion}
 \mu_X([0,t])=W^*(t)-K_X(t)+o(1).
\end{equation}
Consequently, the threshold density exists if and only if the
genuine proper-witness collision surplus $K_X(t)$ converges.
\end{proposition}

\begin{proof}
A represented target absent from the support of $r_t^*$ can only
be represented at threshold $t$ using cofactor $e=1$; such a
target belongs to $\operatorname{range}(\sigma)$. Lemma
\ref{lem:dadd:sigma-range} therefore shows that
\[
 \mu_X([0,t])
 =\frac1X\#\{N\leq X:r_t^*(N)>0\}+o(1).
\]
For every nonnegative integer $r$,
$\boldsymbol1_{\{r>0\}}=r-(r-1)_+$. Sum this identity over
$N\leq X$ with $r=r_t^*(N)$ and apply
Proposition~\ref{prop:dadd:witness-means}.
\end{proof}

\begin{proposition}[A certificate-free positive collision surplus]
\label{prop:dadd:collision-one-ninth}
For every $t\geq2$,
\begin{equation}
\label{eq:dadd:collision-one-ninth}
 \liminf_{X\to\infty}K_X(t)>\frac19.
\end{equation}
\end{proposition}

\begin{proof}
Put $Q=30$, $\Delta=\prod_{p\mid Q}(1-p^{-1})=4/15$, and
\[
 \mathcal C=\{2,4,8,10\}.
\]
For $D\in\mathcal C$ and $(u,Q)=1$, consider the canonical
cofactor-two witness
\begin{equation}
\label{eq:dadd:collision-affine}
 N=s(Du)=\sigma(D)\sigma(u)-Du
        =u\bigl(s(D)+\sigma(D)(h(u)-1)\bigr).
\end{equation}
Its source $m=Du$ is even, so $N=s(m)\geq m/2$ and its
source-to-target ratio is at most $2\leq t$.

Let $\nu_Q$ be the unnormalized Davenport law restricted to
$(u,Q)=1$, of total mass $\Delta$. Expanding $h(u)$ as a divisor
sum gives
\[
 \frac1\Delta\int(h-1)\,d\nu_Q(h)
 =\sum_{\substack{d\geq1\\(d,Q)=1}}\frac1{d^2}-1
 =\frac{8\pi^2}{75}-1=:H.
\]
The passage to this first moment is justified, for example, by
the uniformly bounded second moment of $h(u)$. The radial
partition argument from Proposition~\ref{prop:dadd:witness-means}
gives the source intensity
\begin{align*}
 W_D
 &:=\lim_{X\to\infty}\frac1X
       \#\{u:(u,Q)=1,\ s(Du)\leq X\}\\
 &=\int\frac{d\nu_Q(h)}{s(D)+\sigma(D)(h-1)}
 \geq\frac{\Delta}{s(D)+\sigma(D)H},
\end{align*}
where the last inequality is Jensen's inequality.

Choose the fixed modulus $q=2^4\cdot3\cdot5^2$.
Lemma~\ref{lem:qadd:fixed-modulus-normality} implies that, apart
from $o(X)$ sources $m\leq2X$, one has
\[
 v_p(\sigma(m))>v_p(D)=v_p(m)
 \qquad(p=2,3,5).
\]
For every remaining target in~\eqref{eq:dadd:collision-affine},
\[
 v_p(N)=v_p(\sigma(m)-m)=v_p(D)
 \qquad(p=2,3,5).
\]
Thus, up to $o(X)$ exceptions, its image is contained in the
cylinder
\[
 \mathcal T_D=\{Dv:(v,30)=1\},
 \qquad d(\mathcal T_D)=\frac\Delta D.
\]
If $r_D(N)$ denotes the number of parameters $u$ giving $N$ in
\eqref{eq:dadd:collision-affine}, it follows that
\[
 \liminf_{X\to\infty}\frac1X
       \sum_{N\leq X}(r_D(N)-1)_+
 \geq\Delta\left(
      \frac1{s(D)+\sigma(D)H}-\frac1D\right).
\]
The source cores $D$ are distinct, so their witnessing pairs are
distinct. The elementary inequality
\[
 \sum_{D\in\mathcal C}(r_D(N)-1)_+
 \leq\left(\sum_{D\in\mathcal C}r_D(N)-1\right)_+
 \leq(r_t^*(N)-1)_+
\]
therefore allows the four lower bounds to be added.

Finally, $\pi<22/7$ gives $H<197/3675$. Since
\[
 (D,s(D),\sigma(D))
 \in\{(2,1,3),(4,3,7),(8,7,15),(10,8,18)\},
\]
all four terms are positive and exact rational arithmetic gives
\begin{align*}
 \liminf_{X\to\infty}K_X(t)
 &>
 \frac4{15}
 \sum_{(D,s,B)}
     \left(\frac1{s+197B/3675}-\frac1D\right)\\
 &=\frac{3492203725711}{31001587618275}
  >\frac19.
\end{align*}
This lower bound proves that the collision term in
\eqref{eq:dadd:collision-criterion} is substantial; it does not
prove that the term converges.
\end{proof}

\begin{remark}[What remains open]
None of the preceding results proves the bounded-cofactor
condition~\eqref{eq:dadd:tightness-criterion}, convergence or
uniqueness of the empirical laws, or full support of a possible
limiting distribution. In particular, the strict collision lower
bound~\eqref{eq:dadd:collision-one-ninth} gives no convergence
information about the collision surplus.
\end{remark}
